\documentclass[11pt,a4paper]{article}
\usepackage[spanish]{babel}
\usepackage{lipsum}
\usepackage{nccmath}
\usepackage{amsmath,amssymb}
\title{Mengukur periode ayunan}
\author{Rudi}
\date{\today}

\begin{document}
\maketitle

\section{Pendahuluan}
    \subsection{Pertanyaan penelitian}
        \subsubsection{Asumsi model}
        Berapa lama satu getaran berlangsung? 
\section{Metode}
    Baris sumber pertama\\
    masih menjadi paragraf yang sama.
    
    Ini \textbf{paragraf} baru. Nilai galat adalah 2\%.
    Tabel~1 tidak memisahkan kata dan nomor.
    
    Simbol khusus: \& \_ \# \$ \{ \}
    
    
    Seteleah tanda \% tidak akan terbaca. % ini tidak terbaca

\section{Metode}
    \begin{itemize}
        \item Tali, beban, dan penyangga
        \item Mistar dan pencatat waktu
    \end{itemize}
    \begin{enumerate}
        \item Ukur panjang dari titik gantung ke pusat beban.
        \item Catat waktu untuk sepuluh getaran.
    \end{enumerate}
    \begin{description}
        \item[Periode] Waktu untuk satu getaran.
        \item[Panjang] Jarak titik gantung ke pusat beban.
    \end{description}
    Data ini\footnote{check} bersifat ilustratif.\footnote{Bukan hasil eksperimen.}
    Periksa pengukuran awal.\marginpar{Ukur ulang}
\section{Hasil dan pembahasan}
    \label{sec:hasil}
    Satu periode adalah $T=t/N$, dengan $N$ jumlah getaran.
    \begin{equation}\label{eq:periode}
    T \approx 2\pi\sqrt{\frac{L}{g}}.
    \end{equation}

    \begin{fleqn}
        \begin{align}
        \int f(x) \, dx = F(x) + C
        \end{align}
    \end{fleqn}

    % \begin{freqn}
    %     \begin{align}
    %         \frac{dF(x)}{dx} = f(x)
    %     \end{align}
    % \end{freqn}
            
    Persamaan~\ref{eq:periode} dan persamaan~\eqref{eq:periode} pada subbab \ref{sec:hasil} berlaku untuk sudut kecil.
\section{Kesimpulan}
    % \noindent \lipsum[5]
    \begin{itemize}
        \item List satu
        \item List dua 
    \end{itemize}
    
    % \lipsum[4]
    Ini di halaman kedua\footnote{mencoba footnote.}.
    \[
        T_{10}^2,\quad \frac{a+b}{c},\quad \binom{n}{k},
        \qquad \left(\frac{T_2}{T_1}\right)^2
    \]
    \[
        ( \int_a^b \frac{\left[f(x)+g(x)\right]}{2}\ dx )^3 = F(x_{2})
    \]
    \[
        \left( \int_a^b \frac{\left[f(x)+g(x)\right]}{2}\ dx \right)^3 = F(x_{2})
    \]

    \[
        A = \begin{pmatrix} 1 & 2 \\ 3 & 4 \end{pmatrix},
        \qquad
        \mathbf{t} = \begin{bmatrix} T_1 \\ T_2 \end{bmatrix} 
        \mathbf{t} = \begin{Bmatrix} T_1 \\ T_2 \end{Bmatrix} 
        \qquad 
        % \mathbf{t} = \begin{cmatrix} T_1 \\ T_2 \end{cmatrix}
    \]



\end{document}